\documentclass[12pt]{article}
\usepackage[utf8]{inputenc}
\usepackage[T1]{fontenc}
\usepackage[margin=1in]{geometry}
\usepackage{setspace}
\usepackage{xcolor}
\usepackage[colorlinks=true, urlcolor=blue!60!black, linkcolor=blue!60!black]{hyperref}
\setlength{\parindent}{0em}
\setlength{\parskip}{0.6em}
\renewcommand{\baselinestretch}{1.0}
\begin{document}
\begin{center}
{\LARGE \textit{BMW of North America, Inc. v. Gore}, 517 U.S. 559 (1996)\par}
\vspace{0.5em}
{\large Case Brief --- Two-Page Edition\par}
\end{center}

\section{Memory Jogger}
The Court for the first time held a state punitive award ``grossly excessive'' under due process---\$2 million on \$4,000 in harm for an undisclosed paint repair---and set three guideposts: reprehensibility, ratio, and comparable sanctions (\href{https://www.courtlistener.com/opinion/118026/bmw-of-north-america-inc-v-gore/}{517 U.S. 559}).

\section{Detailed Case Facts}
Dr. Ira Gore bought a new BMW in Alabama and later learned it had been repainted before sale. Under a nationwide policy, BMW sold cars as new without disclosure when repairs cost under 3\% of MSRP; it had sold 983 such cars nationwide, 14 in Alabama. A dealer valued the loss at \$4,000, and Gore asked for \$4 million in punitive damages by multiplying that loss by about 1,000 cars. BMW's threshold complied with roughly 25 state disclosure statutes.

\section{Procedural History}
The jury awarded \$4,000 in compensatory and \$4 million in punitive damages. The Alabama Supreme Court found the nationwide multiplier improper and remitted the award to \$2 million. The Supreme Court granted certiorari and reversed and remanded.

\section{Judicial Votes}
Reversed, 5--4. Stevens wrote for O'Connor, Kennedy, Souter, and Breyer; Breyer, joined by O'Connor and Souter, concurred; Scalia (with Thomas) and Ginsburg (with Rehnquist) dissented.

\section{Holding}
The \$2 million award violated due process. A State may not punish lawful out-of-state conduct, and a defendant must have fair notice of the conduct punished and the severity of the penalty, measured by reprehensibility, the ratio of harm to award, and comparable civil or criminal penalties.

\section{Analysis of Opinions}
Stevens confined Alabama to protecting its own consumers. Reprehensibility, the most important indicium, was low: purely economic harm, no safety risk, a reasonable reading of state statutes as safe harbors, and no deliberate falsehood. The ratio was 500:1, and Alabama's maximum civil penalty was \$2,000. Breyer found that Alabama's standards failed to constrain discretion. Scalia called the guideposts ``a road to nowhere'' and rejected substantive review of awards; Ginsburg would have left reform to the States, which were already enacting caps.

\section{Examples: Future Applications}
\textbf{Same-Side:} a \$15 million punitive award for a \$10,000 insurance underpayment, a 1,500:1 ratio, with a \$5,000 statutory fine, is excessive. \textbf{Opposite-Side:} a 4:1 punitive award against a drugmaker that hid fatal side effects and falsified FDA filings survives. \textbf{Fence-Sitter:} \$150,000 for sabotage causing \$500 in actual but \$200,000 in potential harm.

\section{Critique}
Originalists see a revival of \textit{Lochner}-style substantive due process that federalized state tort law in the name of federalism. Progressives argue that single-plaintiff ratios underdeter widespread low-dollar fraud and that the vague guideposts produced unpredictable results until \textit{State Farm v. Campbell} imposed a presumptive single-digit ceiling.

\section{Key Quotations}
\begin{quote}
``Elementary notions of fairness enshrined in our constitutional jurisprudence dictate that a person receive fair notice not only of the conduct that will subject him to punishment, but also of the severity of the penalty that a State may impose.''
\end{quote}
\noindent --- Justice Stevens, Opinion of the Court, 517 U.S. at 574.

\end{document}
